% !TeX root = ../../金融经济学_王江_讲义及习题.tex
% 来源：教材印刷页237—239，PDF物理页252—254；保留六题题号与原始数据。
% 15.1补全显式标度解；15.2补无套利区间、严格正财富与财富补偿的标量求根式。
% 债券统一按到期支付1、价格B=1/R计数，修正旧解将筹资金额误称债券数量的口径。
% 15.4的alpha比较静态依赖投入尺度；15.5按原题股票支付总产出处理，并补明(c)中参与者2的来源。
% 15.5区分固定股价与固定状态价格，并用原数值反例检验发行量受限时的内点候选。
% 复核原书15.6：恢复题干的“证明”要求；15.5(a)明确驻点与全局最优的求解顺序。
% 15.6外部债券持有H与发行D以D-H进入预算，修正原解购买债券的现金流符号。

\Needspace{10\baselineskip}
\par\noindent\begin{minipage}{\linewidth}
\begin{problem}
\WJFifteenExercise{exer:wj-fifteen-1}{} 考虑一个具有两个时期 0 和 1 的经济。在 1 期有两个可能状态 $a$ 和 $b$，发生概率分别为 $\pi_{a}$ 和 $\pi_{b}$。存在生产技术如下：如果在 0 期投入 $I$，那么在 1 期的产出将是
\[
    Y=\begin{cases}
        \lambda_{a}I^{\alpha},&\text{状态 }a\\
        \lambda_{b}I^{\alpha},&\text{状态 }b
    \end{cases}
\]
其中 $0<\lambda_{b}<\lambda_{a}<1$ 且 $0<\alpha<1$。假设参与者 1 拥有这项技术但却没有任何其它资产。他具有如下形式的效用函数：
\[
    \log(c_{0})+\pi_{a}\log c_{1a}+\pi_{b}\log c_{1b}.
\]

证券市场上只存在无风险债券，它的利率为 $r_{F}$。
\begin{enumerate}[label=(\alph*)]
    \item 写出参与者 1 的优化问题。
    \item 求解他的最优生产、消费和投资选择。
    \item 说明他的生产决策如何依赖于利率。并给出解释。
    \item 计算在最优选择下他的期望效用。
\end{enumerate}
\end{problem}
\end{minipage}\par
\Needspace{5\baselineskip}
\begin{solution}

记 $R=1+r_F>0$，$\pi_a,\pi_b>0$ 且和为1。
\begin{enumerate}[label=(\alph*)]
\item 设发行债券的到期偿付额为 $D$。参与者没有当期禀赋，0 期筹资 $D/R$ 要同时支付当期消费和生产投入。1 期各状态都要偿还 $D$，故求解
\[
 \begin{aligned}
 \max_{I>0,c_0>0,D}\quad &\log c_0+\sum_\omega\pi_\omega\log c_\omega\\
 \text{s.t.}\quad &
 \left\{\begin{aligned}
 c_0&=D/R-I,\\
 c_\omega&=\lambda_\omega I^\alpha-D>0.
 \end{aligned}\right.
 \end{aligned}
\]

由当期预算解得 $D=R(c_0+I)$，代入未来消费，得到 $c_\omega=\lambda_\omega I^\alpha-R(c_0+I)$。这样可直接选择 $c_0$ 和 $I$，借款额由两者决定。
\item 令 $q=D/I^\alpha$，表示每单位 $I^\alpha$ 对应的到期偿付额。因为 $D=R(c_0+I)>0$，且低产出状态的消费必须为正，所以 $0<q<\lambda_b$。各状态消费为 $I^\alpha(\lambda_\omega-q)$，记
\[
 H(q)=\frac{\pi_a}{\lambda_a-q}+\frac{\pi_b}{\lambda_b-q}.
\]

按 (a) 消去 $D$ 后，对 $c_0$ 求导，得 $1/c_0=R\sum_\omega\pi_\omega/c_\omega$。对 $I$ 求导，则有
\[
 \sum_\omega\frac{\pi_\omega}{c_\omega}
       \bigl(\alpha\lambda_\omega I^{\alpha-1}-R\bigr)=0
 \qquad\text{（生产的一阶条件）}.
\]
把消费写成 $I^\alpha(\lambda_\omega-q)$，并用消费一阶条件替换生产条件中的 $R\sum_\omega\pi_\omega/c_\omega$，得到
\[
 \frac1{c_0}=R I^{-\alpha}H(q),\qquad
 \frac1{c_0}=\sum_\omega\frac{\pi_\omega\alpha\lambda_\omega I^{\alpha-1}}{c_\omega}
 =\frac\alpha I[1+qH(q)].
\]

第一式给出 $Rc_0/I^\alpha=1/H(q)$，所以预算 $R(c_0+I)=qI^\alpha$ 化为 $I/c_0=qH(q)-1$。第二条一阶条件则给出 $I/c_0=\alpha[1+qH(q)]$。两式相等，得到
\[
 qH(q)=\frac{1+\alpha}{1-\alpha}.
\]

左侧在 $0<q<\lambda_b$ 上从0严格增至无穷，故有唯一根 $q_0$。记 $C=(1+\alpha)/(1-\alpha)$、$L=\pi_a\lambda_b+\pi_b\lambda_a$，它也是二次方程
\[
 (C+1)q^2-[C(\lambda_a+\lambda_b)+L]q+C\lambda_a\lambda_b=0
\]

落在 $(0,\lambda_b)$ 内的根。最优选择为
\[
 \boxed{I_1^*=\left(\frac{2\alpha q_0}{(1+\alpha)R}\right)^{1/(1-\alpha)}},\qquad
 c_{1,0}^*=\frac{1-\alpha}{2\alpha}I_1^*,\quad
 c_{1,\omega}^*=(I_1^*)^\alpha(\lambda_\omega-q_0).
\]

债券头寸为 $-D_1^*=-q_0(I_1^*)^\alpha$，即卖出相应到期支付。消去 $D$ 后，每个 $c_\omega=\lambda_\omega I^\alpha-R(c_0+I)$ 都是凹函数，正消费可行域为凸集。对 $c_0$ 的对数项及产出对 $I$ 的严格凹性使目标严格凹，因此上述可行的一阶解是唯一最优解。
\item $q_0$ 与利率无关，故
\[
 \frac{d\log I_1^*}{d\log R}=-\frac1{1-\alpha}<0.
\]

利率上升提高融资的偿付成本，最优生产规模下降。
\item 将最优消费代入，得
\[
 U_1^*=\log\frac{1-\alpha}{2\alpha}
 +(1+\alpha)\log I_1^*
 +\sum_\omega\pi_\omega\log(\lambda_\omega-q_0).
\]

这给出由题目参数直接计算的最优期望效用。\qedhere
\end{enumerate}
\end{solution}

\Needspace{10\baselineskip}
\par\noindent\begin{minipage}{\linewidth}
\begin{problem}
\WJFifteenExercise{exer:wj-fifteen-2}{} 在练习~\ref{exer:wj-fifteen-1}描述的经济中，现在假设证券市场还包括一只证券，它在状态 $a$ 支付 $\lambda_{a}$、在状态 $b$ 支付 $\lambda_{b}$ 而它的价格为 $S$。
\begin{enumerate}[label=(\alph*)]
    \item 证明现在证券市场是完全的。
    \item 用 $r_{F}$ 和 $S$ 表示状态价格。
    \item 描述参与者 1 的优化问题。
    \item 计算他的总财富。证明它总是正的。
    \item 求解他的最优生产、消费和投资选择。
    \item 用 NPV 法则解释他的投资选择。
    \item 他的生产技术的 NPV 是多少？它如何随着状态价格的改变而改变？解释所得到的结论。
    \item 计算引入风险证券参与者 1 得到的期望效用的增加值，用 0 期财富的确定性等价表示。
\end{enumerate}
\end{problem}
\end{minipage}\par
\Needspace{5\baselineskip}
\begin{solution}

沿用 $R=1+r_F$，令 $B=1/R$。
\begin{enumerate}[label=(\alph*)]
\item 两证券的支付矩阵为 $\left(\begin{smallmatrix}1&\lambda_a\\1&\lambda_b\end{smallmatrix}\right)$，行列式为 $\lambda_b-\lambda_a\ne0$。由行满秩判据\ifdefined\WJHandoutKnowledge（第~\ref{sec:wj-three-rank}~节）\fi，任意两状态支付都可复制，市场完全。
\item 状态价格由 $\phi_a+\phi_b=B$、$\lambda_a\phi_a+\lambda_b\phi_b=S$ 决定：
\[
 \phi_a=\frac{RS-\lambda_b}{R(\lambda_a-\lambda_b)},\qquad
 \phi_b=\frac{\lambda_a-RS}{R(\lambda_a-\lambda_b)}.
\]

无套利要求两个状态价格都严格为正。由于 $R(\lambda_a-\lambda_b)>0$，这等价于两个分子都为正，即 $\lambda_b<RS<\lambda_a$。以下取这一价格区间。
\item 令债券和风险证券的总持仓为 $\theta_B,\theta_S$，则
\[
 c_0+I+B\theta_B+S\theta_S=0,\qquad
 c_\omega=\lambda_\omega I^\alpha+\theta_B+\lambda_\omega\theta_S.
\]

将未来两条消费式分别乘以 $\phi_a,\phi_b$ 后加总。由 (b)，债券与风险证券支付的现值恰好等于它们的交易成本，即
\[
 \phi_ac_a+\phi_bc_b=SI^\alpha+B\theta_B+S\theta_S .
\]
代入当期预算以消去证券头寸。反过来，因为市场完全，任何满足下式的消费计划都可以用证券复制其相对产出的净支付。因此个人问题为\ifdefined\WJHandoutKnowledge（见式~\eqref{eq:wj-15-4} 的现值预算）\fi
\[
 \begin{aligned}
 \max_{\substack{I\ge0\\c_0,c_a,c_b>0}}\quad
   &\log c_0+\pi_a\log c_a+\pi_b\log c_b\\
 \text{s.t.}\quad &c_0+\phi_ac_a+\phi_bc_b=SI^\alpha-I.
 \end{aligned}
\]
\item 技术创造的最大财富为 $w^*=\max_{I\ge0}(SI^\alpha-I)$。将净值写成 $I^\alpha(S-I^{1-\alpha})$。因为 $S>0$、$1-\alpha>0$，充分小的正投入满足 $I^{1-\alpha}<S$，故最大财富 $w^*>0$。
\item 投入 $I$ 只通过右边的财富影响消费机会。状态价格给定时，财富越高，能够负担的消费越多，所以先选择使 $SI^\alpha-I$ 最大的投入，再分配财富。该净值严格凹，生产条件 $\alpha S(I_2^*)^{\alpha-1}=1$ 给出
\[
 I_2^*=(\alpha S)^{1/(1-\alpha)},\qquad w^*=\frac{1-\alpha}{\alpha}I_2^*.
\]

再在预算 $c_0+\sum_\omega\phi_\omega c_\omega=w^*$ 下选择消费。令预算乘子为 $\nu$，一阶条件为 $1/c_0=\nu$、$\pi_\omega/c_\omega=\nu\phi_\omega$。消去 $\nu$，得到 $\phi_\omega c_\omega=\pi_\omega c_0$。未来消费支出的现值合计等于 $c_0$，因此
\[
 c_0^*=\frac{w^*}{2},\quad c_\omega^*=\frac{\pi_\omega w^*}{2\phi_\omega}.
\]

风险证券和债券头寸为
\[
 \theta_S^*=\frac{c_a^*-c_b^*}{\lambda_a-\lambda_b}-(I_2^*)^\alpha,\qquad
 \theta_B^*=c_a^*-\lambda_a\bigl[(I_2^*)^\alpha+\theta_S^*\bigr].
\]

两式同时复制所需终期消费，0期预算随状态价格条件成立。
\item 由 NPV 法则\ifdefined\WJHandoutKnowledge（定理~\ref{thm:wj-fifteen-1}）\fi，增加一单位投入的边际收益为 $\alpha SI^{\alpha-1}$，边际成本为 $1$。差额为正时增加投入；随着边际产出递减，差额在 $I_2^*$ 处降为零。
\item 最大NPV等于 $w^*$，且包络定理给出
\[
 \frac{\partial w^*}{\partial\phi_\omega}=\lambda_\omega(I_2^*)^\alpha>0.
\]

状态价格上升提高给定产出的现值，重新优化后，投入随 $S=\sum_\omega\phi_\omega\lambda_\omega$ 上升而增加。若固定 $r_F$，则 $\phi_a+\phi_b=B$ 不变。沿着提高 $\phi_a$、等量降低 $\phi_b$ 的方向，最大净现值对 $\phi_a$ 的导数为 $(\lambda_a-\lambda_b)(I_2^*)^\alpha$。
\item 完全市场下的最优效用为
\[
 U_2^*=2\log\frac{w^*}{2}+\sum_\omega\pi_\omega\log\frac{\pi_\omega}{\phi_\omega}.
\]

用 0 期财富衡量开放风险证券带来的效用增益。保持原债券市场的利率，给参与者额外财富 $z$，并允许他重新选择生产、借款和消费。记这样得到的最优效用为 $U_B(z)$，所求补偿 $\Delta w$ 满足 $U_B(\Delta w)=U_2^*$。此时预算为 $c_0=z+D/R-I$。消去 $D$ 后，$z$ 只作为常数进入消费式，因此两条一阶条件与练习~\ref{exer:wj-fifteen-1} 相同。联立 $\frac1{c_0}=RI^{-\alpha}H(q)=\frac{\alpha}{I}[1+qH(q)]$，再代入预算，得
\[
 \begin{aligned}
 I(q)&=\left[\frac{\alpha}{R}\left(q+\frac1{H(q)}\right)\right]^{1/(1-\alpha)},\\
 c_0(q)&=\frac{I(q)^\alpha}{RH(q)},\qquad
 z(q)=I(q)+c_0(q)-\frac{qI(q)^\alpha}{R},\\
 U(q)&=\log c_0(q)+\alpha\log I(q)
 +\sum_\omega\pi_\omega\log(\lambda_\omega-q).
 \end{aligned}
\]

在 $q\le q_0$ 上，令 $M(q)=q+1/H(q)>0$，可将财富写成 $z(q)=I(q)[1+\alpha-2q/M(q)]/\alpha$。由于 $H'>H^2$（严格的柯西--施瓦茨不等式），$M$ 递减。又因 $q/M=qH/(1+qH)$，且 $(qH)'=\sum_\omega\pi_\omega\lambda_\omega/(\lambda_\omega-q)^2>0$，$q/M$ 递增。因此 $z(q)$ 严格下降至 $z(q_0)=0$。当 $q\to-\infty$ 时，$M(q)$ 趋于 $\sum_\omega\pi_\omega\lambda_\omega>0$，故 $z(q)\to\infty$。这给出每个 $z\ge0$ 对应的可行一阶解，严格凹性保证其最优。

新增风险证券扩大可行消费集合，故 $U_2^*\ge U_B(0)=U(q_0)$。$U_B(z)$ 随 $z$ 严格递增，且保持未来消费不变、将新增财富用于当期消费即可使效用无界增长。故可唯一求得 $q\le q_0$ 使 $U(q)=U_2^*$，财富补偿为 $\Delta w=z(q)$。若原最优选择已达到完全市场的最优效用，补偿为零。\qedhere
\end{enumerate}
\end{solution}

\Needspace{10\baselineskip}
\par\noindent\begin{minipage}{\linewidth}
\begin{problem}
\WJFifteenExercise{exer:wj-fifteen-3}{} 考虑练习~\ref{exer:wj-fifteen-1}描述的经济。现在假设经济中除了那里描述的参与者 1 以外还有一个参与者 2。参与者 2 没有生产技术，但经济赋予了他 $e_{0}$ 单位的 0 期消费。他的效用函数和参与者 1 的一样。证券市场和练习~\ref{exer:wj-fifteen-2}中的一样。
\begin{enumerate}[label=(\alph*)]
    \item 求解两个参与者在给定状态价格下的最优生产、消费和投资选择。
    \item 求解均衡状态价格、总消费和消费配置。
    \item 证明均衡配置是 Pareto 最优的。
\end{enumerate}
\end{problem}
\end{minipage}\par
\Needspace{5\baselineskip}
\begin{solution}

设 $e_0>0$，且两状态概率均为正。
\begin{enumerate}[label=(\alph*)]
\item 给定 $\phi_a,\phi_b>0$，记 $S=\sum_\omega\phi_\omega\lambda_\omega$、$B=\sum_\omega\phi_\omega$。沿用练习~\ref{exer:wj-fifteen-2} 的生产净现值最大化，参与者 1 选择 $I(\phi)=(\alpha S)^{1/(1-\alpha)}$，其财富为 $w_1(\phi)=SI(\phi)^\alpha-I(\phi)$。参与者 2 无生产技术，财富为 $w_2=e_0$。按两人的对数消费需求，
\[
 c_{k,0}=\frac{w_k}{2},\qquad c_{k,\omega}=\frac{\pi_\omega w_k}{2\phi_\omega}.
\]

记 $Y_{1,\omega}=\lambda_\omega I(\phi)^\alpha$、$Y_{2,\omega}=0$。用证券复制各人所需的净支付 $x_{k,\omega}=c_{k,\omega}-Y_{k,\omega}$，得到
\[
 \theta_{k,S}=\frac{x_{k,a}-x_{k,b}}{\lambda_a-\lambda_b},\qquad
 \theta_{k,B}=x_{k,a}-\lambda_a\theta_{k,S}.
\]

债券按到期支付 1 计数。两种证券的净供给都为零。
\item 当期资源只有参与者 2 的禀赋 $e_0$，既要供两人消费，也要用于生产投入。因此当期出清为 $c_{1,0}^*+c_{2,0}^*+I^*=e_0$。代入 $c_{1,0}^*=w_1^*/2$、$c_{2,0}^*=e_0/2$，得 $w_1^*=e_0-2I^*$。另一方面，生产条件给出 $w_1^*=(1-\alpha)I^*/\alpha$，所以
\[
 I^*=\frac{\alpha e_0}{1+\alpha},\qquad
 C_0^*=\frac{e_0}{1+\alpha},\qquad
 C_\omega^*=\lambda_\omega(I^*)^\alpha.
\]

未来市场出清与两人的消费需求还给出
\[
 \frac{\pi_\omega(w_1^*+e_0)}{2\phi_\omega^*}=\lambda_\omega(I^*)^\alpha.
\]

两人面对同样的状态价格且偏好相同，(a) 的需求表明，各节点的消费份额都等于财富份额 $w_k^*/(w_1^*+e_0)$。代入 $w_1^*=(1-\alpha)e_0/(1+\alpha)$，得到 $\eta_1=(1-\alpha)/2$、$\eta_2=(1+\alpha)/2$，所以
\[
 c_{k,0}^*=\eta_kC_0^*,\qquad c_{k,\omega}^*=\eta_kC_\omega^*,
 \qquad
 \boxed{\phi_\omega^*=\frac{\pi_\omega e_0^{1-\alpha}}
 {\lambda_\omega\alpha^\alpha(1+\alpha)^{1-\alpha}}}.
\]

由 $B^*=\phi_a^*+\phi_b^*$、$S^*=\lambda_a\phi_a^*+\lambda_b\phi_b^*$ 恢复证券价格与 $r_F^*=1/B^*-1$。代回生产条件可核对 $\alpha S^*(I^*)^{\alpha-1}=1$。各人终期消费与技术产出均与 $(\lambda_a,\lambda_b)$ 成比例，因此 $\theta_{1,B}^*=\theta_{2,B}^*=0$，$\theta_{2,S}^*=\eta_2(I^*)^\alpha=-\theta_{1,S}^*$，两种证券均出清。
\item 用正权重检验均衡配置的帕累托最优性\ifdefined\WJHandoutKnowledge（定理~\ref{thm:wj-fifteen-2}）\fi。取 $\eta_1,\eta_2$ 为规划权重。给定投入 $I$，在任一总消费为 $C$ 的节点，规划者在 $c_1+c_2=C$ 下最大化 $\sum_k\eta_k\log c_k$。一阶条件使 $\eta_1/c_1=\eta_2/c_2$，故 $c_k=\eta_k C$。将最优分享代回，节点效用为 $\log C$ 加一个与 $I$ 无关的常数。再代入当期和未来总资源，生产选择化为
\[
 \max_{0<I<e_0}\ \log(e_0-I)+\alpha\log I+\sum_\omega\pi_\omega\log\lambda_\omega.
\]

它严格凹，一阶条件给出上述 $I^*$；因此该配置最大化正权重的效用和，是帕累托最优的。\qedhere
\end{enumerate}
\end{solution}

\Needspace{10\baselineskip}
\par\noindent\begin{minipage}{\linewidth}
\begin{problem}
\WJFifteenExercise{exer:wj-fifteen-4}{} 考虑练习~\ref{exer:wj-fifteen-3}中的经济。
\begin{enumerate}[label=(\alph*)]
    \item 构建代表性参与者。
    \item 通过求解代表性参与者的优化问题而得到均衡配置。
    \item 用代表性参与者的边际效用计算状态价格。
    \item 讨论 $\lambda_{b}$、$\lambda_{a}$ 以及 $\alpha$ 如何影响状态价格和均衡配置。
\end{enumerate}
\end{problem}
\end{minipage}\par
\Needspace{5\baselineskip}
\begin{solution}

\begin{enumerate}[label=(\alph*)]
\item 沿用上一题得到的规划权重 $\eta_k$ 和分享规则。给定某一期或状态的总消费 $C$，把最优分配 $c_k=\eta_k C$ 代回该节点的加权效用，得到
\[
 \max_{\substack{c_1+c_2=C\\c_1,c_2>0}}\sum_k\eta_k\log c_k
 =\log C+\sum_k\eta_k\log\eta_k.
\]

舍去不影响选择的常数，代表性参与者效用为 $\log C_0+\sum_\omega\pi_\omega\log C_\omega$。
\item 当期总消费 $C_0=e_0-I$ 必须为正，所以 $I<e_0$。未来各状态的总消费 $C_\omega=\lambda_\omega I^\alpha$ 也须为正，因此 $I>0$。把这两条资源关系代入代表性效用，得到
\[
 \begin{aligned}
 \max_I\quad &\log(e_0-I)+\sum_\omega\pi_\omega\log(\lambda_\omega I^\alpha)\\
 \text{s.t.}\quad &0<I<e_0.
 \end{aligned}
\]

一阶条件给出
\[
 -\frac1{e_0-I}+\frac\alpha I=0,
 \qquad I^*=\frac{\alpha e_0}{1+\alpha}.
\]

目标对 $I$ 严格凹，故该内点为唯一最优解。代回资源约束得 $C_0^*=e_0/(1+\alpha)$、$C_\omega^*=\lambda_\omega(I^*)^\alpha$，个人消费为 $c_{k,0}^*=\eta_k C_0^*$、$c_{k,\omega}^*=\eta_k C_\omega^*$。
\item 对代表性参与者的两期消费一阶条件\ifdefined\WJHandoutKnowledge（式~\eqref{eq:wj-15-6}）\fi，消去预算乘子。状态价格就等于未来该状态的边际效用 $\pi_\omega/C_\omega^*$ 除以当期的边际效用 $1/C_0^*$，即
\[
 \phi_\omega^*=\pi_\omega\frac{C_0^*}{C_\omega^*}
 =\frac{\pi_\omega e_0^{1-\alpha}}
 {\lambda_\omega\alpha^\alpha(1+\alpha)^{1-\alpha}}.
\]
\item 提高 $\lambda_\omega$ 按比例增加该状态的产出和消费，并降低其状态价格；$I^*,C_0^*$ 和消费份额保持不变。

提高 $\alpha$ 时，$dI^*/d\alpha=e_0/(1+\alpha)^2>0$，$dC_0^*/d\alpha=-e_0/(1+\alpha)^2<0$；个人份额满足 $d\eta_1/d\alpha=-1/2$、$d\eta_2/d\alpha=1/2$。对 $\log C_\omega^*=\log\lambda_\omega+\alpha\log I^*$ 求导时，要同时考虑指数 $\alpha$ 和投入 $I^*$ 的变化。再用 $\log\phi_\omega^*=\log\pi_\omega+\log C_0^*-\log C_\omega^*$，得到
\[
 \frac{d\log C_\omega^*}{d\alpha}=\log I^*+\frac1{1+\alpha},\qquad
 \frac{d\log\phi_\omega^*}{d\alpha}=-\log I^*-\frac2{1+\alpha}.
\]

个人当期消费中，$c_{2,0}^*=e_0/2$ 保持不变，$c_{1,0}^*=C_0^*-e_0/2$ 随 $\alpha$ 上升而下降。由于 $c_{k,\omega}^*=\eta_k C_\omega^*$，个人未来消费的对数变化率在总消费的对数变化率上，分别加上 $-1/(1-\alpha)$ 和 $1/(1+\alpha)$。

保持生产函数的计量单位不变。若 $I^*>\exp[-1/(1+\alpha)]$，提高 $\alpha$ 会增加未来总消费；若 $I^*>\exp[-2/(1+\alpha)]$，提高 $\alpha$ 会降低状态价格。低于相应阈值时方向相反，恰等于阈值时相应导数为零。\qedhere
\end{enumerate}
\end{solution}

\Needspace{10\baselineskip}
\par\noindent\begin{minipage}{\linewidth}
\begin{problem}
\WJFifteenExercise{exer:wj-fifteen-5}{} 沿用练习~\ref{exer:wj-fifteen-2} 中的生产技术与偏好。在这里我们不假设有一组给定的可交易证券，而是假设拥有生产技术的参与者 1 可以成立一个公司、而公司可以到市场上发行债券和股票，但没有其它可交易的证券。一单位（无风险）公司债在 1 期有 1 单位支付，而一股股票在 1 期有 $\hat{\lambda}I^{\alpha}$ 单位支付，在状态 $a$ 和 $b$ 下 $\hat{\lambda}=\lambda_{a}$，$\lambda_{b}$，$I$ 是公司的生产投资。
\begin{enumerate}[label=(\alph*)]
    \item 给定利率 $r_{F}$ 和股价 $S$，求解参与者 1 的优化问题，即最优的投资额、在市场发行的债券和股票的数量以及他的消费。
    \item 在这种情况下，有无最优负债水平？这和 Modigliani--Miller 定理是否一致？解释原因。
    \item 求解市场均衡，特别地，求解均衡利率、股价、生产水平、发行的债券和股票的数量以及两个参与者的消费配置。
\end{enumerate}
\end{problem}
\end{minipage}\par
\Needspace{5\baselineskip}
\begin{solution}

% 原题股票支付为全部产出lambda I^alpha，非扣债后的剩余支付；按题设沿用，发行者承担债务偿付。
令 $R=1+r_F>0$，$B=1/R$。以 $D$ 表示到期支付1的公司债发行量，$s$ 表示出售的总产出份额。依题设，发行者保留 $(1-s)\lambda_\omega I^\alpha$，并从中偿付债务。
\begin{enumerate}[label=(\alph*)]
\item 给定 $R,S$，当期消费来自债券筹资 $BD$ 和股票筹资 $Ss$ 扣除投入 $I$ 后的余额。未来消费是保留的产出扣除债务偿付后的余额。题目未规定债股头寸的上下限，先按头寸可正可负写出问题，其中 $D<0$ 表示持有债券：
\[
 \begin{aligned}
 \max_{I,D,s}\quad &\log c_{1,0}+\sum_\omega\pi_\omega\log c_{1,\omega}\\
 \text{s.t.}\quad &
 \left\{\begin{aligned}
 c_{1,0}&=BD+Ss-I,\\
 c_{1,\omega}&=(1-s)\lambda_\omega I^\alpha-D,\\
 c_{1,0},c_{1,\omega}&>0\quad(\forall\omega),\\
 I&\ge0.
 \end{aligned}\right.
 \end{aligned}
\]

先按本问字面，对所有投入都固定同一个 $S>0$。投入越小，每股支付越低，但售价不变。取足够小的 $I>0$，使一股股票的售价 $S$ 高于用债券买到其最高状态支付所需的成本 $B\lambda_a I^\alpha$。若头寸不限，可卖出 $t$ 股，并买入 $t\lambda_a I^\alpha+d$ 单位债券，其中 $d>0$。即取 $s=t$、$D=-t\lambda_a I^\alpha-d$，代入两期预算得到
\[
 c_{1,0}=t(S-B\lambda_a I^\alpha)-Bd-I,\qquad
 c_{1,\omega}=\lambda_\omega I^\alpha+d+t(\lambda_a-\lambda_\omega)I^\alpha.
\]

当 $t\to\infty$ 时，当期消费趋于无穷，两状态的未来消费始终至少为 $d$，故不存在有限最优选择。这个构造利用了不限量出售股票并持有债券的假设。若将“发行”理解为 $D\ge0$、$0\le s\le1$，则要把这些限制加入可行域。以下保留正投入、正消费时的内点计算，供核查发行限制下的候选解。

当债股发行限制均不绑定时，一阶条件为
\[
 \frac B{c_{1,0}}=\sum_\omega\frac{\pi_\omega}{c_{1,\omega}},\quad
 \frac S{c_{1,0}}=I^\alpha\sum_\omega\frac{\pi_\omega\lambda_\omega}{c_{1,\omega}},\quad
 \frac1{c_{1,0}}=\alpha(1-s)I^{\alpha-1}\sum_\omega\frac{\pi_\omega\lambda_\omega}{c_{1,\omega}}.
\]

将投入条件与股票条件相除，得 $I=\alpha S(1-s)$。再将股票条件乘以 $1-s$、债券条件乘以 $D$ 后相减，右侧为 $\sum_\omega\pi_\omega c_{1,\omega}/c_{1,\omega}=1$。因此 $S(1-s)-BD=c_{1,0}$。与当期预算相加，得到 $2c_{1,0}=S-I$。于是可将内点条件化为单变量方程
\[
 s(I)=1-\frac I{\alpha S},\quad
 D(I)=\frac R{2\alpha}\bigl[(\alpha+2)I-\alpha S\bigr],\quad
 t(I)=\frac{I^{1+\alpha}}{\alpha S},
\]
\[
 \sum_\omega\frac{\pi_\omega}{t(I)\lambda_\omega-D(I)}
 =\frac2{R(S-I)}.
\]

固定 $I>0$ 时，消费是 $(D,s)$ 的仿射函数，对数效用使融资子问题严格凹。若加上发行限制，需先在每个 $I$ 下求融资子问题的最大值，再对 $I$ 求全局最大值，并检查发行限制绑定的解。上述方程只给出无约束一阶条件的候选。
\item 按固定股价且头寸不限的读法，(a) 没有有限最优解，因此也没有最优债务水平。若限制发行，固定 $I>0$ 时，融资子问题在有解时至多有一个最优债券发行量和出售份额。按 MM 定理\ifdefined\WJHandoutKnowledge（定理~\ref{thm:wj-fifteen-4}）\fi，比较融资安排时，应固定真实投资和总产出。这里改变 $I$ 就会改变每股支付，而报价 $S$ 仍被固定，所以 (a) 的投资问题无法用一组共同状态价格为所有可选支付定价。
\item 原题在此问引入第二位参与者，以下沿用练习~\ref{exer:wj-fifteen-3} 中参与者 2 的禀赋 $e_0$ 与偏好。按固定 $R,S$ 且头寸不限的读法，均衡必须有 $I>0$ 以满足未来正消费，还必须有 $S>0$，否则参与者 2 可不限量买入正支付的股票。由 (a)，在这样的价格下参与者 1 没有最优选择，故不存在该口径下的均衡。

\textbf{发行受限时的均衡检验。}令发行者满足 $D\ge0$、$0\le s\le1$，参与者 2 可按报价自由交易两种证券，包括净供给为零时。记 $Y_\omega=\lambda_\omega I^\alpha$。证券出清时，参与者 2 的消费为
\[
 c_{2,0}=e_0-BD-Ss,\qquad c_{2,\omega}=D+sY_\omega.
\]

未来总消费为 $Y_\omega$，故 $I>0$。$s=1$ 会使发行者未来消费非正，故 $s<1$。若 $s=0$，发行者须以 $D>0$ 筹资，参与者 2 的未来消费在两态都是 $D$。买方对债券和股票的一阶条件给出 $S=B\mathrm E[\widetilde Y]$。此时发行者的债券选择为内点，满足 $B/c_{1,0}=\mathrm E[1/\widetilde c_1]$，其中 $\widetilde c_1$ 为发行者的未来消费。将两式代入发行者对 $s$ 的边际效用，
\[
 \left.\frac{\partial U_1}{\partial s}\right|_{s=0}
 =\frac{S}{c_{1,0}}-\mathrm E[\widetilde Y/\widetilde c_1]
 =-\operatorname{Cov}(\widetilde Y,1/\widetilde c_1)>0.
\]

最后一步是因为 $\widetilde c_1=\widetilde Y-D$，产出越高，消费的倒数越低。因此发行少量股票可提高效用，排除 $s=0$。

若 $D>0$，则发行者对 $D,s$ 的选择均为内点。两只证券的支付线性独立，所以两人的债股一阶条件给出同一组状态价格。两人的对数消费需求使未来消费各按固定比例分享总产出。若参与者 2 的份额为 $\eta_2$，则 $D+sY_\omega=\eta_2Y_\omega$。由于两状态的 $Y_\omega$ 不同，两式相减得 $s=\eta_2$，再代回得 $D=0$，与 $D>0$ 矛盾。

因此只剩 $D=0$、$0<s<1$。两人的未来消费分别为 $(1-s)Y_\omega$ 和 $sY_\omega$，故按个人边际效用计算的两个状态价格向量成比例。双方对股票的一阶条件都将这组价格乘以产出支付后加总为 $S>0$，比例系数只能为 $1$，所以两组状态价格相同。

参与者 2 的现值财富是 $e_0$，按练习~\ref{exer:wj-fifteen-2} 的对数消费需求，当期消费为 $e_0/2$。两人的状态价格相同，也给出相同的未来与当期消费比例，所以总产出的市值等于当期总消费 $e_0-I$。再由发行者对 $I,s$ 的一阶条件，有
\[
 S^{\dagger}s^{\dagger}=\frac{e_0}{2},\qquad
 S^{\dagger}=e_0-I^{\dagger},\qquad I^{\dagger}=\alpha S^{\dagger}(1-s^{\dagger}).
\]

联立得到候选
\[
 \boxed{I^{\dagger}=\frac{\alpha e_0}{2(1+\alpha)},\quad
 S^{\dagger}=\frac{(2+\alpha)e_0}{2(1+\alpha)},\quad
 s^{\dagger}=\frac{1+\alpha}{2+\alpha},\quad D^{\dagger}=0}.
\]
\[
 \begin{aligned}
 c_{1,0}^{\dagger}&=\frac{e_0}{2(1+\alpha)},& c_{2,0}^{\dagger}&=\frac{e_0}{2},\\
 c_{1,\omega}^{\dagger}&=\frac{\lambda_\omega(I^{\dagger})^\alpha}{2+\alpha},&
 c_{2,\omega}^{\dagger}&=\frac{(1+\alpha)\lambda_\omega(I^{\dagger})^\alpha}{2+\alpha}.
 \end{aligned}
\]

总消费 $C_0^{\dagger}=S^{\dagger}$、$C_\omega^{\dagger}=\lambda_\omega(I^{\dagger})^\alpha$ 给出
\[
 \phi_\omega^{\dagger}=\frac{\pi_\omega S^{\dagger}}{\lambda_\omega(I^{\dagger})^\alpha},\qquad
 B^{\dagger}=\sum_\omega\phi_\omega^{\dagger},\qquad r_F^{\dagger}=1/B^{\dagger}-1.
\]

这是固定股价、发行受限设定下唯一可能的均衡候选。还需在这组报价下检验发行者是否全局最优。

例如取 $\alpha=0.6$、$e_0=10$、$(\lambda_a,\lambda_b)=(0.8,0.3)$、$(\pi_a,\pi_b)=(0.4,0.6)$，候选为
\[
 (I,D,s)=(1.875,0,0.6153846),\quad S=8.125,\quad R\simeq0.07178561.
\]

发行者在该价格下的效用约为 $-0.25055307$。改选可行头寸
\[
 (I,D,s)=(1.88355771,\ 0.00286777,\ 0.61107863),
\]

效用约为 $-0.25053986$，高于候选配置的效用。这次改进满足 $D\ge0$、$0\le s\le1$，各期各状态的消费均为正。因此唯一候选未通过最优性检验，\textbf{该组参数下，固定股价的受限模型不存在均衡}。

\textbf{固定状态价格时的均衡。}若交易者接受的是一组给定的状态价格 $\phi$，每个投入下的一股股票都按对应的实际支付定价。只要允许为实现状态消费而作所需交易，股价就随投入变化
\[
 S(I)=I^\alpha\sum_\omega\phi_\omega\lambda_\omega.
\]

此时，最大化净值 $S(I)-I$ 给出与练习~\ref{exer:wj-fifteen-3} 相同的生产条件及状态价格 $\phi_\omega^*=\pi_\omega C_0^*/C_\omega^*$。代入 $C_0^*=e_0-I^*$、$C_\omega^*=\lambda_\omega(I^*)^\alpha$，得到
\[
 I^*=\frac{\alpha e_0}{1+\alpha},\qquad S^*=C_0^*=\frac{e_0}{1+\alpha},\qquad
 r_F^*=\frac{1}{\phi_a^*+\phi_b^*}-1 .
\]

个人消费为 $c_{k,0}^*=\eta_k C_0^*$、$c_{k,\omega}^*=\eta_k C_\omega^*$，其中 $\eta_1=(1-\alpha)/2$、$\eta_2=(1+\alpha)/2$。取 $D^*=0$、$s^*=\eta_2$，参与者 2 买入全部售出的股票，付款 $S^*s^*=e_0/2$，取得 $\eta_2 C_\omega^*$。参与者 1 取得 $S^*s^*-I^*=\eta_1 C_0^*$ 并保留 $\eta_1 C_\omega^*$。这样债股交易均出清，消费加总后恰好等于各期资源。\qedhere
\end{enumerate}
\end{solution}

\Needspace{28\baselineskip}
\par\noindent\begin{minipage}{\linewidth}
\begin{problem}
\WJFifteenExercise{exer:wj-fifteen-6}{} 在练习~\ref{exer:wj-fifteen-5}中，现在假设市场上存在一只无风险证券可供参与者交易。证明此时Modigliani--Miller 定理成立，即公司的融资决策不影响它的投资决策和均衡配置。解释这种情形和上述情形之间的差别。
\end{problem}
\end{minipage}\par
\Needspace{5\baselineskip}
\begin{solution}

令 $D$ 为公司债到期偿付额，$H$ 为另行购买的外部无风险证券数量。两种债券每单位到期都支付 $1$，由一价定律\ifdefined\WJHandoutKnowledge（定理~\ref{thm:wj-four-3}）\fi，价格同为 $B=1/(1+r_F)$。公司发债带来当期收入 $BD$ 和未来偿付 $D$，个人买入外部债券则带来当期支出 $BH$ 和未来收入 $H$。因此个人预算为
\[
 c_{1,0}=B\underbrace{(D-H)}_{\text{净负债}}+Ss-I,\qquad
 c_{1,\omega}=(1-s)\lambda_\omega I^\alpha-(D-H).
\]

给定投入 $I$ 和出售份额 $s$，两种债券对消费的影响由净负债 $L=D-H$ 决定。上一题的净负债就是公司债发行量 $D$。引入外部债券后，只要调整可行，同时增加 $D$ 和 $H$ 同一数量就能维持原来的消费，因而改变公司债发行量也能保留相同的投资与个人净头寸。

这一步说明给定 $I$ 时，外部债券怎样抵销公司融资。要进一步比较不同的 $I$，还要为随投入变化的股票支付定价。若仍对所有投入固定同一股价 $S$，取 $H=0$ 就能沿用上一题的效用改进。

沿用上一题的状态定价口径，令 $B=\sum_\omega\phi_\omega$、$S(I)=I^\alpha\sum_\omega\phi_\omega\lambda_\omega$。将未来消费按状态价格加总，现值为 $(1-s)S(I)-B(D-H)$。加上当期预算后，净债务的收入与偿付现值相消，售出和保留的产出价值合为 $S(I)$，于是
\[
 c_{1,0}+\sum_\omega\phi_\omega c_{1,\omega}
 =\underbrace{S(I)-I}_{\text{生产创造的净值}}.
\]

给定正投入，债券与总产出证券张成两个状态。只要个人债券交易可以抵销公司融资的变动，各种融资安排对应同一个预算集合，最优投入都由最大化 $S(I)-I$ 决定。

例如，沿用上一题固定状态价格下的均衡，设新增外部债券的净供给为零，且允许借贷。公司增发 $\Delta$ 单位债券，参与者 1 同时增持 $\Delta$ 单位外部债券。参与者 2 买入全部增发的公司债，同时卖空 $\Delta$ 单位外部债券。两人各自在 0 期和 1 期的新增现金流均相抵，公司债和外部债券也都出清。只要这组头寸可行，投资、股价和两人的均衡消费保持上一题的配置，公司的债务发行量可以改变。\qedhere
\end{solution}

